\documentclass[11pt,a4paper]{article}
\usepackage[margin=24mm]{geometry}
\usepackage{amsmath,amssymb,booktabs,tabularx,microtype}
\usepackage[T1]{fontenc}
\usepackage{lmodern}
\usepackage[colorlinks=true,urlcolor=blue,linkcolor=blue]{hyperref}
\setlength{\parindent}{0pt}
\setlength{\parskip}{6pt}
\title{What is the wave in the black-hole scattering project?}
\author{Project note}
\date{}
\begin{document}
\maketitle
\vspace{-1.5em}

\section{The field we evolve}

The current project evolves a \textbf{massless scalar field on a fixed
Schwarzschild spacetime}. The wave is a disturbance in the scalar field.
A small perturbation of the metric would instead describe the related
problem of gravitational perturbations.

In units \(G=c=1\), the background metric is
\begin{equation}
 ds^2=-f(r)\,dt^2+f(r)^{-1}\,dr^2+r^2\,d\Omega^2,
 \qquad f(r)=1-\frac{2M}{r}.
\end{equation}
We keep this metric fixed and solve
\begin{equation}
 \Box_{g_{\mathrm{Schw}}}\varphi
 =\frac{1}{\sqrt{-g}}\partial_\mu
 \left(\sqrt{-g}\,g^{\mu\nu}\partial_\nu\varphi\right)=0.
\end{equation}
This is a controlled test-field approximation: if
\(\Phi=\varepsilon\varphi\), then the scalar stress-energy is quadratic
in \(\varepsilon\). With no independent gravitational perturbation,
\begin{equation}
 T_{\mu\nu}[\Phi]=O(\varepsilon^2),
 \qquad g_{\mu\nu}=g^{\mathrm{Schw}}_{\mu\nu}+O(\varepsilon^2).
\end{equation}
Thus the scalar can propagate on the background at leading order while
its effect on the geometry is neglected. A unit code amplitude is a
normalisation of this linear problem.

\section{How the three-dimensional wave becomes a radial problem}

Separate the angular dependence using a spherical harmonic:
\begin{equation}
 \varphi(t,r,\theta,\phi)=\frac{\psi_{\ell m}(t,r)}{r}
 Y_{\ell m}(\theta,\phi).
\end{equation}
For a real scalar, use real spherical harmonics or a real combination of
complex harmonics. Each mode evolves independently. Define the
\emph{tortoise coordinate}
\begin{equation}
 x=r_*=r+2M\ln\left(\frac{r}{2M}-1\right).
\end{equation}
The exterior \(2M<r<\infty\) maps to \(-\infty<x<\infty\).
Suppressing the mode indices, the equation becomes
\begin{equation}
 \boxed{\quad \partial_t^2\psi=\partial_x^2\psi-V_\ell(r(x))\psi,\qquad
 V_\ell(r)=f(r)\left[\frac{\ell(\ell+1)}{r^2}
                         +\frac{2M}{r^3}\right].\quad}
\end{equation}
This is a one-dimensional wave equation with a spatially varying
potential. The potential accounts for angular structure and curvature;
it produces scattering outside the horizon. It is not a material
surface. For Schwarzschild, the radial potential depends on \(\ell\)
but not on \(m\).

\clearpage
\section{Specifying the wave means specifying initial data}

A second-order equation in time requires \textbf{two initial functions}:
\begin{equation}
 \psi(0,x)=\psi_0(x),\qquad
 \partial_t\psi(0,x)=\Pi_0(x).
\end{equation}
The code's two evolved fields will be
\(\psi\) and \(\Pi=\partial_t\psi\).
A convenient initial radial profile is a Gaussian:
\begin{equation}
 \psi_0(x)=A\exp\left[-\frac{(x-x_0)^2}{2\sigma^2}\right].
\end{equation}
Here \(A\) is the amplitude, \(x_0\) the centre, and \(\sigma\) the width.
The shape alone does not determine which way the wave travels.

Where the potential is negligible, the general solution is
\begin{equation}
 \psi(t,x)=F(x-t)+G(x+t).
\end{equation}
The first term moves towards increasing \(x\); the second moves towards
decreasing \(x\), hence towards the black hole. An initially inward-moving
Gaussian therefore uses
\begin{equation}
 \boxed{\quad \Pi_0(x)=+\partial_x\psi_0(x)
              =-\frac{x-x_0}{\sigma^2}\psi_0(x).\quad}
\end{equation}
The minus sign choice, \(\Pi_0=-\partial_x\psi_0\), selects outward motion.
Choosing \(\Pi_0=0\) gives equal inward and outward parts in the free-wave
limit. At finite radius, a nonzero potential means these are local
free-wave prescriptions, not exact pure scattering states.

A single angular mode is a radial pulse with an angular pattern:
\(\ell=0\) is spherically symmetric, while \(\ell=2\) has quadrupolar
structure. A wave packet localised on just one side of the black hole
would require a superposition of angular modes.

\section{Properties you can vary}

\begin{tabularx}{\textwidth}{@{}lX@{}}
\toprule
Parameter & Effect in the experiment \\
\midrule
\(A\) & Scales the field linearly and its energy quadratically; does not
change the background's ringdown frequencies. \\
\(\sigma\) & Changes the pulse width and frequency content. A narrower
Gaussian has a broader frequency spectrum. \\
\(x_0\) & Changes the launch position, arrival time and initial overlap
with the scattering potential. \\
\(\Pi_0\) & Selects initial inward motion, outward motion, or a mixture. \\
\(\ell\) & Changes the angular structure, potential and ringdown spectrum. \\
\(k_0\) & Adds an oscillation scale to a Gaussian wave packet. \\
\bottomrule
\end{tabularx}

For the last option, use
\begin{equation}
 \psi_0(x)=A e^{-(x-x_0)^2/(2\sigma^2)}
             \cos\!\left[k_0(x-x_0)\right],
 \qquad \Pi_0=\partial_x\psi_0.
\end{equation}
Far from the potential, \(k_0\) sets the approximate carrier frequency.
Varying it lets you study how the reflected and transmitted energy
fractions depend on frequency.

\clearpage
\section{How to see the propagation}

Once time evolution is implemented, three complementary views will show
what the field does:
\begin{enumerate}
 \item \textbf{Snapshots or an animation of \(\psi(t,x)\).}
 Watch the initial pulse move inward, interact with the potential, and
 split into inward and outward components.
 \item \textbf{A spacetime heat map of \(\psi(t,x)\).}
 Put \(x\) on the horizontal axis, \(t\) on the vertical axis, and field
 value in colour. In regions where the potential is small, propagation
 follows \(x-t=\mathrm{constant}\) or \(x+t=\mathrm{constant}\).
 \item \textbf{A waveform at a fixed observer.}
 Record
 \begin{equation}
   w(t)=\psi(t,x_{\mathrm{obs}}).
 \end{equation}
 An observer outside the barrier can see the outgoing scattered signal
 and its subsequent decaying oscillations, called \emph{ringdown}.
\end{enumerate}

The one-dimensional plots show the reduced radial field. To reconstruct
the scalar at a spacetime point, restore the angular factor and \(1/r\):
\begin{equation}
 \varphi(t,r,\theta,\phi)=\frac{\psi(t,r_*(r))}{r}Y_{\ell m}(\theta,\phi).
\end{equation}
The waveform's time coordinate is Schwarzschild \(t\); a static observer
at radius \(r_{\mathrm{obs}}\) has proper time
\(d\tau=\sqrt{f(r_{\mathrm{obs}})}\,dt\).

The transmitted part moves towards \(x\to-\infty\), corresponding to the
horizon. In the planned finite numerical domain, the first experiment
uses distant reflecting endpoints and a time window chosen to exclude
significant returning boundary reflections. A larger-domain comparison
will check that the measured scattering is independent of those walls.

\section{Connection to gravitational waves and the project goal}

If instead we perturb the metric,
\begin{equation}
 g_{\mu\nu}=g^{\mathrm{Schw}}_{\mu\nu}
             +\varepsilon h_{\mu\nu}+O(\varepsilon^2),
\end{equation}
the leading vacuum problem is the linearised Einstein equation
\(\delta G_{\mu\nu}[h]=0\). Schwarzschild gravitational perturbations can
be described by Regge--Wheeler and Zerilli master equations of the form
\begin{equation}
 \partial_t^2\Psi=\partial_x^2\Psi-V_{\mathrm{grav}}(r)\Psi.
\end{equation}
The numerical structure is similar, but the variables and potentials
have different physical meanings. This offers a natural later extension;
the present scalar waveform is not a gravitational-wave strain.

\textbf{The current goal} is to build and validate a C++/AMReX solver for
scalar scattering, extract its ringdown, and study numerical accuracy
and adaptive refinement. The first practical exercise remains a flat,
periodic travelling wave with \(V=0\), before introducing Schwarzschild
curvature.

\smallskip
\small
Further reading: E.~Berti, V.~Cardoso and A.~O.~Starinets,
\emph{Quasinormal modes of black holes and black branes},
\href{https://arxiv.org/abs/0905.2975}{arXiv:0905.2975}.
The repository's longer theory guide contains the detailed derivations.
\end{document}
